% ============================================================
% ABSTRACT
% ============================================================
\begin{abstract}
% Test-time compute improves decision quality for LLM agents, but deciding when it is beneficial remains unresolved. Existing gating methods rely on uncertainty signals, implicitly assuming a fixed relationship between signal value and compute benefit. We show that this assumption is incorrect. 
% Across six diverse agent environments and three model backbones, we find that the same signal value that predicts rollout benefit in one environment predicts rollout harm in another, indicating a reversal in the signal--utility relationship. 
% We trace this to a conflation of two uncertainty sources, \emph{information poverty} and \emph{decision difficulty}, and formalize it via a two-source model grounded in Simpson's paradox, where subgroup-level trends reverse upon aggregation. We show that any gating strategy with a fixed signal–utility direction must underperform never-gating in at least one such environment, making direction discovery unavoidable for cross-environment robustness.
% Motivated by this result, we propose DIAL (Direction-Informed Adaptive Learning), which learns environment-specific signal--utility relationships through direction discovery. Across six environments, DIAL outperforms all fixed-direction baselines with zero additional per-step overhead, achieving adaptive gating behavior solely by learning signal direction. 

% Test-time compute improves decision quality for LLM agents, but deciding \emph{when} to invoke it across diverse environments remains an open problem. Existing adaptive-gating methods share an implicit assumption: that higher uncertainty consistently indicates when extra computation helps. We show this assumption fails in heterogeneous agent settings. Across six environments and three model backbones, the same signal value that predicts rollout benefit in one environment predicts rollout harm in another, and once misaligned, better calibration \emph{worsens} performance by more precisely selecting harmful states.

% We trace this reversal to two coexisting state types: \emph{information-poor} states, where the agent lacks evidence and rollouts cannot compensate, and \emph{decision-difficult} states, where multiple viable options exist and rollouts help explore them. These produce opposite signal--utility relationships, so any aggregate correlation depends on their mixture (an instance of Simpson's paradox). Since environments differ in this mixture, no gate with a fixed signal--utility direction can avoid falling below the no-trigger baseline in at least one environment, making per-environment direction discovery necessary.

% Motivated by this, we propose DIAL (Direction-Informed Adaptive Learning), which learns environment-specific signal--utility relationships from interaction data. Across six environments, DIAL outperforms all fixed-direction baselines while adding zero per-step overhead at deployment.

% Test-time compute improves decision quality for LLM agents, but deciding \emph{when} to invoke it across diverse environments remains an open problem. Existing adaptive-gating methods share an implicit assumption: that higher uncertainty consistently indicates when extra computation helps. We show this assumption fails in heterogeneous agent settings: the same signal value that predicts rollout benefit in one environment predicts rollout harm in another, and once misaligned, better calibration \emph{worsens} performance by more precisely selecting harmful states.
% We trace this reversal to two coexisting state types: \emph{information-poor} states, where the agent lacks evidence and rollouts cannot compensate, and \emph{decision-difficult} states, where multiple viable options exist and rollouts help explore them. These produce opposite signal--utility relationships, so any aggregate correlation depends on their mixture (an instance of Simpson's paradox). Since environments differ in this mixture, no gate with a fixed signal--utility direction can avoid falling below the no-trigger baseline in at least one environment, making per-environment direction discovery necessary.
% Motivated by this, we propose DIAL (Direction-Informed Adaptive Learning), which learns environment-specific signal--utility relationships from interaction data. DIAL outperforms all fixed-direction baselines while adding zero per-step overhead at deployment.
For LLM agents, \emph{better-calibrated uncertainty signals can make adaptive test-time compute strictly worse than no gating at all}. Existing adaptive-gating methods share an implicit fixed-direction assumption---that higher uncertainty consistently indicates when extra computation helps---and we show this assumption fails in heterogeneous agent settings: the same signal value that predicts rollout benefit in one environment predicts rollout harm in another, and the sign can flip across model backbones with the task held fixed. When the assumed direction is wrong, improving calibration \emph{worsens} performance by more precisely selecting harmful states---the gate becomes a more accurate harm-targeting mechanism.
We trace this reversal to two coexisting state types with opposite signal--utility relationships---\emph{information-poor} states, where rollouts cannot compensate for missing evidence, and \emph{decision-difficult} states, where rollouts help explore viable alternatives---whose mixture is determined jointly by the environment and the backbone. The decomposition yields a sharper, information-theoretic claim: \emph{the uncertainty signal is not a sufficient statistic for the value of computation}, so no policy depending on uncertainty alone---monotone or otherwise, calibrated or not---can deliver non-negative value across heterogeneous (environment, backbone) settings. Recovering sufficiency requires augmenting the signal with auxiliary features that distinguish the two state types.
As an existence proof that this augmentation is the bottleneck---not gate sophistication---we instantiate the minimal procedure consistent with the analysis: DIAL (Direction-Informed Adaptive Learning), a sparse linear gate fit on signal-agnostic exploration data over a small pool of universal features. Across six environments and three backbones, DIAL already outperforms every fixed-direction baseline at zero per-step inference overhead, with its full range of gating behavior---from aggressive triggering to near-total suppression---emerging from direction learning alone.
\end{abstract}
